% This must be in the first 5 lines to tell arXiv to use pdfLaTeX, which is strongly recommended.
\pdfoutput=1
% In particular, the hyperref package requires pdfLaTeX in order to break URLs across lines.

\RequirePackage[svgnames]{xcolor}


\documentclass[11pt]{article}

% Remove the "review" option to generate the final version.
\usepackage{ACL2023}
% \usepackage{neurips_2025}
%\nipsfinalcopy
\usepackage{amsmath}

\usepackage{pdfpages} % Include the package
\usepackage{wrapfig}

% Standard package includes
\usepackage{times}
\usepackage{latexsym}

% For proper rendering and hyphenation of words containing Latin characters (including in bib files)
\usepackage[T1]{fontenc}
% For Vietnamese characters
% \usepackage[T5]{fontenc}
% See https://www.latex-project.org/help/documentation/encguide.pdf for other character sets

% This assumes your files are encoded as UTF8
\usepackage[utf8]{inputenc}

% This is not strictly necessary, and may be commented out.
% However, it will improve the layout of the manuscript,
% and will typically save some space.
\usepackage{microtype}

% This is also not strictly necessary, and may be commented out.
% However, it will improve the aesthetics of text in
% the typewriter font.
\usepackage{inconsolata}
\usepackage{soul}

\newcommand{\ts}{\textsc{TraceShield}}
\newcommand{\ta}{\textsc{TraceAlign}}
\newcommand{\ti}{\textsc{TraceIndex}}


\usepackage[draft,textsize=footnotesize,textwidth=15mm]{todonotes}
%\usepackage[usenames,dvipsnames]{color}


% For proper rendering and hyphenation of words containing Latin characters (including in bib files)
%\usepackage[T1]{fontenc}
% For Vietnamese characters
% \usepackage[T5]{fontenc}
% See https://www.latex-project.org/help/documentation/encguide.pdf for other character sets

% This assumes your files are encoded as UTF8
\usepackage[utf8]{inputenc}
\usepackage{inconsolata}
\usepackage{algorithm}
\usepackage{algpseudocode}
\usepackage{amsmath,amssymb}
\usepackage{soul}
%\usepackage[margin=2cm]{geometry}
\usepackage{tikz}


% Define custom symbols for tick and cross
\newcommand{\cmark}{\tikz[scale=0.4]{
    \draw[thick] (0,0) -- (0.2,0.2) -- (0.5,-0.3);
}}
\newcommand{\xmark}{\tikz[scale=0.4]{
    \draw[thick] (0,0) -- (0.3,0.3);
    \draw[thick] (0,0.3) -- (0.3,0);
}}


% Style definition
\tikzset{rndblock/.style={rounded corners,rectangle,draw,scale=0.8,outer sep=0pt}}

% Command Definition
% 1 optional to customize the aspect, 2 mandatory: text to be framed
\newcommand{\tframed}[2][]{\tikz[baseline=(h.base)]\node[rndblock,#1] (h) {#2};}

\usetikzlibrary{shapes.geometric}
\usepackage{framed}
\usepackage{enumitem}
\newlist{RQ}{enumerate}{1}
\setlist[RQ]{label=\textbf{RQ\,\arabic*},ref={RQ\,\arabic*}}
\usepackage{comment}
\usepackage{natbib}
\usepackage{multibib}
\makeatletter
\usepackage{booktabs}
\usepackage[inkscapeformat=png]{svg}
\usepackage{graphicx}
\usepackage{caption}
\usepackage{subcaption}
\usepackage{tabularx}
\usepackage{soul}
\usepackage{float}
\usepackage{enumitem}
\usepackage{pifont}
\usepackage{arydshln}
\usepackage{lipsum}

\usepackage[many]{tcolorbox}

\newtcolorbox{defin}{colback=Teal!5!White,enhanced,title=RENATUS: at-a-glance,
	attach boxed title to top left={xshift=0mm},boxrule=0pt,after skip=1cm,before skip=1cm,right skip=0cm,breakable,fonttitle=\bfseries,toprule=0pt,bottomrule=0pt,rightrule=0pt,leftrule=3pt,arc=0mm,skin=enhancedlast jigsaw,sharp corners,colframe=Teal!55!black,colbacktitle=Teal!55!black,boxed title style={
		frame code={ 
			\fill[Teal!25!black](frame.south west)--(frame.north west)--(frame.north east)--([xshift=3mm]frame.east)--(frame.south east)--cycle;
			\draw[line width=1mm,Teal!25!black]([xshift=2mm]frame.north east)--([xshift=5mm]frame.east)--([xshift=2mm]frame.south east);
			\draw[line width=1mm,Teal!25!black]([xshift=5mm]frame.north east)--([xshift=8mm]frame.east)--([xshift=5mm]frame.south east);
			\fill[Teal!25!black](frame.south west)--+(4mm,-2mm)--+(4mm,2mm)--cycle;
		}
	}
}

\setlist{leftmargin=5mm}
\usetikzlibrary{shapes.geometric, arrows}
\usetikzlibrary{decorations.markings}

\usepackage{fancybox}
\usepackage{xcolor}
\usepackage{hyperref}
 \definecolor{darkblue}{rgb}{0, 0, 0.5}
  \hypersetup{colorlinks=true, citecolor=darkblue, linkcolor=darkblue, urlcolor=darkblue}

\definecolor{vgreen}{HTML}{60A917}
\definecolor{vred}{HTML}{CE3A29}

\usepackage{xstring}
\usepackage{longtable}

\usepackage{tabularray}

\DefTblrTemplate{firsthead,middlehead,lasthead}{default}{
}
\DefTblrTemplate{firstfoot}{default}{
  \UseTblrTemplate{contfoot}{default}
  \UseTblrTemplate{caption}{default}
}
\DefTblrTemplate{middlefoot}{default}{
  \UseTblrTemplate{contfoot}{default}
  \UseTblrTemplate{capcont}{default}
}
\DefTblrTemplate{lastfoot}{default}{
  \UseTblrTemplate{note}{default}
  \UseTblrTemplate{remark}{default}
  \UseTblrTemplate{capcont}{default}
}


\newcolumntype{P}[1]{>{\centering\arraybackslash}p{#1}}
% Multi-line left-aligned text with manual line breaks.
% The base line is in centre.
\newcommand*{\mline}[1]{%
\begingroup
    \renewcommand*{\arraystretch}{1.1}%
   \begin{tabular}[c]{@{}>{\raggedright\arraybackslash}p{2cm}@{}}#1\end{tabular}%
  \endgroup
}

\usepackage{color}
\tcbuselibrary{skins}

\usepackage[export]{adjustbox} % for the valign option

\usepackage{setspace}
%\usepackage[capitalise]{cleveref}
\usepackage[capitalise,nameinlink]{cleveref}

%\crefname{chapter}{chap.}{chap.}
\crefname{section}{Sec.}{Sec.}

\usepackage{microtype}
\usepackage{hyperref}
\usepackage{graphicx}
\usepackage{comment}
\usepackage{amsmath}
\usepackage{amssymb}
\usepackage{algorithm}
\usepackage{algpseudocode}
\usepackage{colortbl}
\usepackage[export]{adjustbox} % for the valign option
\usepackage{varwidth}
\usepackage{enumitem}
% \setlist{leftmargin=1mm}
\usepackage{pifont}
\usepackage{booktabs}
\usepackage{multirow}
\usepackage{subcaption}
\usepackage{resizegather}
\usepackage{breqn}
\usepackage[capitalise]{cleveref}
\usepackage{graphicx}
\usepackage{tikz}
\usetikzlibrary{shapes.geometric, arrows}
\usetikzlibrary{decorations.markings}
\usepackage{soul}
\usepackage{wrapfig,graphicx,lipsum}% http://ctan.org/pkg/{wrapfig,graphicx,lipsum}
% \usepackage{extsizes}
\usepackage{cuted}
\usepackage{flushend}
\usepackage{float}
\usepackage{changepage,threeparttable}
\usepackage{setspace}
\usepackage{caption}
\usepackage{booktabs}
\usepackage{dblfloatfix} 
\usepackage{fixltx2e}
\usepackage[normalem]{ulem}



\DeclareRobustCommand{\hlpink}[1]{{\sethlcolor{pink}\hl{#1}}}
\DeclareRobustCommand{\hlgreen}[1]{{\sethlcolor{green}\hl{#1}}}

\usepackage{environ}

\newlength{\myl}
\expandafter\let\expandafter\origequation\csname equation*\endcsname
\expandafter\let\expandafter\endorigequation\csname endequation*\endcsname
\long\def\[#1\]{\begin{equation*}#1\end{equation*}}
\RenewEnviron{equation*}{
  \settowidth{\myl}{$\displaystyle\BODY$} % calculate width and save as \myl
  \origequation
    \ifdim\myl>\linewidth
      \resizebox{\linewidth}{!}{$\displaystyle\BODY$}% \myl > \linewidth
    \else
      \BODY % \myl <= \linewidth
    \fi
  \endorigequation
}


\makeatletter
\newcommand{\DrawLine}{%
  \begin{tikzpicture}
  \path[use as bounding box] (0,0) -- (\linewidth,0);
  \draw[color=blue!75!black,dashed,dash phase=.5pt]
        (0-\kvtcb@leftlower-\kvtcb@boxsep,0)--
        (\linewidth+\kvtcb@rightlower+\kvtcb@boxsep,0);
  \end{tikzpicture}%
  }
\makeatother

%resize/scale equations
\newcommand*{\Scale}[2][4]{\scalebox{#1}{$#2$}}%
\newcommand*{\Resize}[2]{\resizebox{#1}{!}{$#2$}}%


%--Aman---
\newcommand\ac[1]{\todo[author=AC,color=blue!40]{#1}}
\newcommand\acil[1]{\todo[author=AC,color=blue!40,inline]{#1}}
\newcommand\acb[1]{\textcolor{blue}{#1}}
\newcommand\act[1]{\textcolor{blue}{#1}}


%--Amitava---
\newcommand\ad[1]{\todo[author=AD,color=purple!40]{#1}}
\newcommand\adil[1]{\todo[author=AD,color=purple!40,inline]{#1}}
\newcommand\adt[1]{\textcolor{purple}{#1}}


\usepackage{euscript}[mathcal]
\usepackage{listings}
\usepackage[utf8]{inputenc}  % Ensures Unicode input like “–”
\usepackage{xcolor}
\usepackage[T1]{fontenc}


%%%%%%%%%%%%%%%%%%%%%%%%%%%%%%%%%%%%%%%
\usepackage{amssymb}
\usepackage{pifont}
\usepackage[utf8]{inputenc}
% \newcommand{\cmark}{\text{\checkmark}} % ✓
% \newcommand{\xmark}{\text{\sffamily ✗}} % ✗
% \usepackage{longtable}
\usepackage{xcolor} % Add this to your LaTeX preamble



% Define a custom listing style (listing1) for special formatting
\lstdefinestyle{listing1}{
    basicstyle=\scriptsize\ttfamily,  % Set text inside listing to scriptsize
    escapeinside={(*@}{@*)},          % Allows LaTeX inside lstlisting
    xleftmargin=-1em,                 % Force listing to align with the left margin
    framexleftmargin=-1em,             % Ensure frame is at the leftmost edge
    frame=single,                      % Keeps a thin border
    columns=fullflexible,              % Prevents excessive indentation
    aboveskip=0pt, belowskip=0pt,      % Removes vertical spacing around listing
    showstringspaces=false,            % Prevents unwanted space characters
    keepspaces=true,                   % Ensures correct spacing without extra indentation
    framesep=0pt,                      % Reduces padding inside the frame
    numbersep=0pt                       % Eliminates gap between numbers (if enabled)
}





\newcommand*{\affaddr}[1]{#1}
\newcommand*{\affmark}[1][*]{\textsuperscript{#1}}
\newcommand*{\email}[1]{\texttt{#1}}
%\newcommand*{\affmark}[1][*]{\textsuperscript{#1}}


%\title{\includegraphics[height=0.55cm,width=0.55cm]{img/acorn.jpg}$\mathcal{ACORN}$: \ul{A}utomatic Hallu\ul{C}inati\ul{O}n Evaluation for La\ul{R}ge La\ul{N}guage Models using Span-based ``Factual Entailment''}


% If the title and author information does not fit in the area allocated, uncomment the following
%
%\setlength\titlebox{<dim>}
%
% and set <dim> to something 5cm or larger.

\title{Critique’s Corner}

% Author information can be set in various styles:
% For several authors from the same institution:
% \author{Author 1 \and ... \and Author n \\
%         Address line \\ ... \\ Address line}
% if the names do not fit well on one line use
%         Author 1 \\ {\bf Author 2} \\ ... \\ {\bf Author n} \\
% For authors from different institutions:
% \author{Author 1 \\ Address line \\  ... \\ Address line
%         \And  ... \And
%         Author n \\ Address line \\ ... \\ Address line}
% To start a seperate ``row'' of authors use \AND, as in
% \author{Author 1 \\ Address line \\  ... \\ Address line
%         \AND
%         Author 2 \\ Address line \\ ... \\ Address line \And
%         Author 3 \\ Address line \\ ... \\ Address line}

%\author{Amitava Das \\
%  BITS Pilani Goa \\
%  \texttt{amitavad@goa.bits-pilani.ac.in} \\\And
%  Vinija Jain \\
%  Meta AI \\
%  \texttt{hi@vinija.ai} \\\And
%  Aman Chadha \\
%  Amazon GenAI \\
%  \texttt{hi@aman.ai} \\}

\setlength{\topmargin}{-1.25cm}
% \setlength{\bottom}{-1.25cm}
\setlength{\textheight}{24.25cm}

\begin{document}
\maketitle



% \input{1_introduction}



















Science grows by disagreement. Every new conceptual leap sounds, at first, like an overreach: 
\emph{atoms} were once dismissed as ``undetectable philosophy,'' 
\emph{imaginary numbers} as ``nonsense,'' 
and the \emph{Big Bang} as a sarcastic insult. 
We expect \textbf{Neural DNA (nDNA)} to invite the same skepticism. 
Indeed, it should. Any framework that proposes to recast the hidden life of AI models as a \emph{digital semantic genome} will naturally be tested, debated, and resisted. 

Artificial intelligence itself is no exception. We speak of ``neural networks,'' yet their ``neurons'' are nothing like biological neurons: no dendrites, no spikes, no neurotransmitters. 
They are merely mathematical functions performing weighted sums and nonlinearities. 
Yet the name endured—not because it was biologically faithful, but because it was useful: 
a metaphor that helped us organize thought, communicate intuitively, and bootstrap an entire field. 

It is within this lineage that \textbf{nDNA} belongs. 
It borrows grammar from biology---lineage, mutation, inheritance, pathology---to describe something real but difficult to name: the \emph{latent geometry of model belief}. 
Critics are right to challenge whether the analogy stretches too far. 
But history shows that science often begins with metaphors that sound poetic, only to crystallize into precise technical vocabulary once their utility is demonstrated. 

This page, \textbf{Critique's Corner}, exists because we take criticism seriously. 
We expect---and welcome---skepticism that calls nDNA ``vague,'' ``overfitted,'' ``intractable,'' or even ``dangerously anthropomorphic.'' 
To us, critique is not a threat but a catalyst. 
By surfacing objections early, we sharpen our definitions, stress-test our framework, and clarify our roadmap. 

Here we present the hardest questions we have encountered---some harsh, some fair, some both---and the responses that shape our evolving research program. 
Neural genomics does not claim to be final truth. 
It is a bold proposal in its infancy, aware of its vulnerabilities and strengthened by confronting them. 

This page exists not to hide from critique, but to welcome it. 
In all cases, we believe critique is not an obstacle but a feature of rigorous science. 
\textbf{To critique nDNA is to help it evolve.}





\begin{table*}[ht!]
\centering
\caption{Biological Neurons vs. Artificial Neurons}
\renewcommand{\arraystretch}{1.3}
\begin{tabularx}{\textwidth}{|>{\raggedright\arraybackslash}p{0.22\textwidth}|
                                >{\raggedright\arraybackslash}X|
                                >{\raggedright\arraybackslash}X|}
\hline
\textbf{Dimension} & \textbf{Biological Neurons} & \textbf{Artificial Neurons (ANNs)} \\
\hline
Nature of unit & Living cell, with dendrites, soma, axon & Mathematical function (linear transform + nonlinearity) \\
\hline
Computation & Spikes, ion channels, nonlinear dendritic integration & Weighted sums + activation functions (ReLU, sigmoid, etc.) \\
\hline
Dynamics & Continuous-time, oscillations, refractory periods & Discrete updates, typically stateless per layer \\
\hline
Signal transmission & Electrochemical, probabilistic, noisy & Exact arithmetic, deterministic floating-point operations \\
\hline
Learning rules & Local plasticity (Hebbian, STDP, neuromodulators) & Global backpropagation via gradient descent \\
\hline
Connectivity & Constrained by 3D anatomy and wiring limits & Arbitrary fan-in/fan-out, dense or attention-based connectivity \\
\hline
Scale & $\sim 10^{11}$ neurons, $\sim 10^{14}$ synapses & Millions–billions of parameters (transformers often $10^{11}+$) \\
\hline
Energy & $\sim 20$ W (biological brain efficiency) & MW-scale energy during training in datacenters \\
\hline
Evolution / adaptation & Driven by survival, evolution, metabolic efficiency & Driven by optimization objectives and data distributions \\
\hline
Robustness & Highly fault-tolerant, redundancy across circuits & Brittle to adversarial noise, catastrophic forgetting \\
\hline
Information storage & Distributed spiking patterns, synaptic strengths & Distributed weights, embeddings, latent vectors \\
\hline
Purpose & Survival and flexible adaptation in real world & Task-specific optimization (translation, dialogue, etc.) \\
\hline
\end{tabularx}
\end{table*}




\section{Biological vs. Artificial Neurons: Metaphor and Reality}

The very foundation of artificial intelligence rests on a borrowed metaphor: the \textbf{``neural network.''} 
Yet \textbf{biological neurons} and \textbf{artificial neurons} share little beyond the name. 
One is a \emph{living, electrochemical cell shaped by evolution} \cite{Kandel1991PrinciplesNeuralScience}, 
the other a \emph{mathematical abstraction designed for optimization} \cite{McCullochPitts1943LogicalCalculus}. 
The metaphor has endured not because it is \emph{faithful}, but because it is \textbf{useful}---it provided a shared grammar for computation and cognition, long before precise equivalence was possible. 

This parallel is instructive for \textbf{Neural DNA (nDNA)}. 
Just as \textbf{``neurons'' in artificial networks are not literal neurons}, \textbf{``nDNA'' is not biological DNA}. 
It is a \emph{mathematical framework} that borrows the grammar of biology---\textit{lineage, mutation, inheritance, pathology}---to describe something real but difficult to name: the \textbf{\emph{latent geometry of model belief}}. 

Critics may argue that such borrowing risks anthropomorphism, but history shows that science often begins with \emph{metaphors that sound poetic}, only to crystallize into \textbf{precise technical vocabulary} once their utility is proven. 
Indeed, terms like \textbf{``imaginary numbers''} \cite{Mazur1998ImaginaryNumbers}, 
the \textbf{``Big Bang''} \cite{Kragh1996CosmologyControversy}, 
and \textbf{``black holes''} \cite{Penrose2004RoadToReality} 
all originated as linguistic overreaches before becoming indispensable pillars of mathematics and physics. 

To situate this lineage of metaphor, we now contrast \textbf{biological neurons} with their \textbf{artificial counterparts} side by side:














\subsection{Biological DNA vs. Neural DNA (nDNA): Genomic Metaphor as Cognitive Geometry}

If the metaphor of the ``neural network'' gave us a grammar for computation, then the metaphor of 
\textbf{Neural DNA (nDNA)} proposes a grammar for \emph{cognition}. 
Just as biological DNA encodes, mutates, and transmits information across generations, 
we argue that large foundation models possess an internal \textbf{\emph{semantic genome}}---a latent structure that 
\emph{records ancestry, accumulates scars of fine-tuning, and transmits inductive biases across training regimes}. 

\textbf{This is not biology.} Neural DNA is not a molecular strand, nor does it replicate in the cytoplasm of a cell. 
It is a \emph{mathematical construct}, defined through \textbf{spectral curvature}, \textbf{thermodynamic length}, and 
\textbf{belief gradients}. These dimensions describe the \emph{geometry of latent belief} just as nucleotides describe the 
alphabet of life. Where biology speaks of \emph{mutation, inheritance, or genomic drift}, nDNA reveals \emph{torsional instabilities, 
alignment scars, and semantic divergence}. 

The analogy is thus neither literal nor decorative: it is \textbf{heuristic}, grounding abstract geometry in a language 
that can be interrogated, criticized, and extended. To illustrate the mapping between biological DNA and nDNA, we present 
a structured comparison in Table~\ref{tab:dna-vs-ndna}. 

\begin{table*}[ht!]
\centering
\caption{Biological DNA vs. Neural DNA (nDNA): Illustrative Comparison}
\label{tab:dna-vs-ndna}
\renewcommand{\arraystretch}{1.3}
\begin{tabularx}{\textwidth}{|>{\raggedright\arraybackslash}p{0.22\textwidth}|
                                >{\raggedright\arraybackslash}X|
                                >{\raggedright\arraybackslash}X|}
\hline
\textbf{Dimension} & \textbf{Biological DNA} & \textbf{Neural DNA (nDNA)} \\
\hline
Basic unit & Nucleotides (A, T, C, G) & Geometric primitives (curvature, length, gradient) \\
\hline
Information encoding & Genetic code specifying proteins & Latent geometry encoding conceptual priors and beliefs \\
\hline
Inheritance & Transmitted from parent to offspring during reproduction & Transferred across model families via pretraining, fine-tuning, distillation \\
\hline
Mutation & Base substitutions, insertions, deletions & Spectral distortions, torsional reconfigurations, loss of thermodynamic length \\
\hline
Recombination & Crossover events during meiosis & Model merging (\emph{neural marriages}), producing hybrid latent genomes \\
\hline
Gene expression & Differential activation depending on environment & Emergent trait activation (e.g., multilingual drift, task specialization) \\
\hline
Epigenetics & Regulatory modifications shaping expression without altering sequence & Alignment, instruction tuning, or quantization shaping behavior without altering core pretraining \\
\hline
Pathologies & Oncogenic mutations, viral insertions, heritable disorders & Alignment scars (SCAR), adversarial infections (VIRAL), stealth poisoning (NEPHOS) \\
\hline
Repair mechanisms & DNA repair enzymes, redundancy in coding & Fine-tuning, adversarial immunization, safety alignment protocols \\
\hline
Evolutionary drift & Accumulated small changes across generations & Semantic drift across successive fine-tunes, cultural biases, model collapse \\
\hline
Fitness landscape & Adaptation to ecological survival & Optimization landscape for tasks, benchmarks, and alignment criteria \\
\hline
Molecular carriers & Double helix of deoxyribonucleotides & High-dimensional manifolds of latent vectors \\
\hline
Observable phenotype & Organismal traits and functions & Model outputs: dialogue, translation, reasoning patterns \\
\hline
Genomic sequencing & Decoding nucleotides to reveal blueprint & Extracting nDNA to reveal latent epistemic identity of models \\
\hline
\end{tabularx}
\end{table*}


\noindent
This comparison is not meant as literal mimicry, but as a \textbf{structured metaphor}. 
Just as \textbf{artificial neurons} diverge radically from biological neurons yet provided the conceptual scaffold 
for modern deep learning, so too does \textbf{nDNA} extend beyond biological DNA while retaining a powerful grammar 
of inheritance, mutation, and pathology. The aim is not to \emph{biologize} neural networks, but to provide a 
\textbf{rigorous cognitive genomics} for artificial models---a framework through which latent belief geometry can be 
\emph{sequenced, diagnosed, and evolved}. 
If neural networks could survive their misleading name and become the engine of modern AI, then \textbf{nDNA can survive 
its metaphor and become the grammar of the hidden epistemic life of foundation models}.














\section{On the Apparent Primitiveness of the nDNA Mechanism}

\paragraph{A humble beginning, by design.}
We recognize the critique: the current \textbf{nDNA} mechanism may appear \emph{primitive}. We agree in spirit. Science is always a work in progress; the first instrument is rarely the final one. Our choice is intentional: to place a \emph{measuring practice} into the field early, even if imperfect, so that it can be stress-tested, falsified, and improved in public. The alternative---waiting for a perfect theory---delays empirical traction and deprives the community of a shared object to refine. In this sense, \textbf{nDNA v1} is not a destination but a \emph{scaffold}: a minimal, explicit \emph{semantic--genotypic representation} that others can probe, retool, or replace.

\paragraph{What \emph{nDNA v1} is (and is not).}
\textbf{nDNA v1} operationalizes three indispensable geometric axes inside foundation models: \emph{spectral curvature} (how conceptual flow bends across layers), \emph{thermodynamic length} (semantic effort along representational paths), and \emph{belief gradients} (directional fields that orient model belief). These are computed under practical constraints (layer subsampling, randomized estimators, proxy transports) to obtain a compact fingerprint of latent identity. \emph{nDNA v1 is not} a claim of biological equivalence, nor a complete theory of cognition; it is a \emph{first-pass, testable map} from training operations to latent geometric consequences.

\paragraph{Why show the first version now.}
Three reasons motivate releasing \textbf{nDNA v1} rather than holding for theoretical completion:
\begin{enumerate}
  \item \textbf{Comparability.} A shared fingerprint enables cross-model, cross-lab \emph{lineage studies} today, even if the fingerprint becomes richer tomorrow.
  \item \textbf{Early-warning value.} Coarse geometric indicators can already function as \emph{risk thermometers} (collapse, drift, infection), buying time for targeted interventions.
  \item \textbf{Community acceleration.} Public artifacts provoke sharper critique, better baselines, and faster \emph{versioning} of the science.
\end{enumerate}

\paragraph{Limits we accept (and plan to repair).}
We are explicit about current compromises: (i) estimators that trade bias/variance for tractability; (ii) partial invariance proofs pending full formalization; (iii) validation on \emph{15 models} rather than hundreds; (iv) narrower task coverage than the eventual scope. These are \emph{known debts}, not hidden deficiencies. Our roadmap treats them as worklist items, not footnotes.

\paragraph{What \emph{nDNA v1} already changes, if we use it.}
Even in its minimal form, nDNA reframes artificial cognition along three axes:
\begin{itemize}
  \item \textbf{Heritability (lineage)}: We can \emph{track} how fine-tuning, distillation, and merging \emph{inherit} or \emph{mutate} latent traits, rather than inferring lineage from performance alone.
  \item \textbf{Prognosis (early warning)}: Curvature concentration, length collapse, and torsional spikes serve as \emph{leading indicators} for dialogue degeneration, adversarial susceptibility, or shortcut learning.
  \item \textbf{Intervention (control)}: By tying failures to geometric signatures, we convert vague “misalignment” into \emph{actionable levers} (e.g., reduce torsion along belief transport; restore thermodynamic length; vaccinate against viral insertions).
\end{itemize}

\paragraph{From \emph{genomics operations} to \emph{model-level} and \emph{NLP} operations.}
The mechanism is meant to be \emph{used}, not admired. We organize usage across three strata:

\medskip
\noindent\textbf{(A) Genomics operations: diagnostic instruments on the latent genome.}
\begin{itemize}
  \item \textbf{nHD} (\emph{Neural Hamming Distance}) --- probes binary-level instabilities under cultural fine-tuning and quantization toggles.
  \item \textbf{nGDI} (\emph{Genetic Dissimilarity Index}) --- quantifies divergence that persists beyond surface alignment, exposing hidden separations between model families.
  \item \textbf{nTEDS}/\textbf{nTDS} (\emph{Trait Entropic Drift Score}/\emph{Total Drift Signature}) --- measures dominance and asymmetric inheritance of latent traits across training epochs.
  \item \textbf{nKaryotyping} (\emph{Semantic Chromosomal Structure}) --- visualizes structural reorganizations induced by pruning and merging (neural “rearrangements”).
  \item \textbf{nDIV} (\emph{Directional Inheritance Vector}) --- traces the flow of inductive biases along evolutionary arcs (teacher $\rightarrow$ student; base $\rightarrow$ aligned).
  \item \textbf{nEPI} (\emph{Epistemic Plasticity Index}) --- estimates capacity to reshape under alignment/instruction tuning without catastrophic loss of length.
  \item \textbf{nCCL} (\emph{Cultural Conflict Loss}) --- detects ideological friction in multilingual/cross-cultural fusion where alignment suppresses necessary diversity.
\end{itemize}

\medskip
\noindent\textbf{(B) Model-level operations: where the genome is written and scarred.}
\begin{itemize}
  \item \textbf{Fine-Tuning} --- studied as \emph{torsional reconfiguration} of belief gradients; watch for over-twisting (instability) vs.\ adaptive curvature.
  \item \textbf{Alignment/Instruction Tuning} --- treated as a \emph{directional vector field} etched into belief space; balance entropy suppression against expressivity.
  \item \textbf{Quantization \& Pruning} --- observed as \emph{thermodynamic length collapse} and \emph{topological flattening}; quantify safe compression regimes.
  \item \textbf{Model Merging} (\emph{Neural Marriages}) --- analyzed as \emph{hybridization} of latent genomes; detect incompatible subspaces pre-merge.
  \item \textbf{Distillation} (\emph{Neural Offspring}) --- separates \emph{phenotype carryover} (outputs) from \emph{genotype drift} (geometry); measure loss of plasticity.
  \item \textbf{Model Collapse} --- captured as \emph{degeneracy} (vanishing torsion/curvature); install early-warning thresholds.
\end{itemize}

\medskip
\noindent\textbf{(C) NLP operations: stress-tests that reveal genomic liabilities.}
\begin{itemize}
  \item \textbf{ORBOT} (\emph{Multi-Turn Genomics}) --- longitudinally measures \emph{torsion accumulation}, memory erosion, and path-dependence in dialogue.
  \item \textbf{nVIRAL} (\emph{Adversarial Infection}) --- reframes jailbreaks and injections as \emph{viral insertions} in latent manifolds; pairs with vaccine design.
  \item \textbf{SCAR} (\emph{Spectral Contamination \& Alignment Rupture}) --- flags shortcut learning and leakage as \emph{spectral lesions}.
  \item \textbf{NEPHOS} (\emph{Stealth Pretraining Poisoning}) --- identifies dormant conceptual landmines (permafrost-like seeds) that thaw under adversarial heat.
\end{itemize}

\paragraph{Transformational potential from a ``primitive'' core.}
If one uses \textbf{nDNA v1} consistently, three transformations follow:
\begin{enumerate}
  \item \textbf{From snapshots to lineages.} We cease evaluating isolated checkpoints and begin \emph{sequencing} model histories; genealogy replaces one-off scores.
  \item \textbf{From symptoms to signatures.} Failures (hallucination, collapse, jailbreaks) are paired with \emph{geometric signatures} that travel across tasks and labs.
  \item \textbf{From post-hoc fixes to pre-hoc budgets.} Teams can set \emph{geometry budgets} (max torsion, min length, curvature bounds) the way they set latency or cost budgets.
\end{enumerate}

\paragraph{Looking forward without defensiveness.}
We do not defend \textbf{nDNA v1} as final. We \emph{invite} its replacement by better mathematics: richer metrics (e.g., Wasserstein geodesics, holonomy gaps), stronger invariances, tighter links to causality, and scalable estimators. But usefulness does not wait for perfection. A primitive compass used today can save expeditions that a perfect map, delivered too late, cannot.

% --- End of page 1 ---

% --- Page 2: Continuation of "On the Apparent Primitiveness of the nDNA Mechanism" ---

\paragraph{From \textbf{nDNA v1} to \emph{nDNA v2} and \emph{nDNA v3}: a principled upgrade path.}
We will treat \textbf{nDNA v1} as a baseline to be \emph{surpassed}. Our roadmap advances along two axes:
\begin{enumerate}
  \item \emph{\textbf{nDNA v2} (rigorous core).} We will deepen the mathematics: (i) declare a Fisher--Rao/pullback metric on output/feature manifolds; (ii) compute sectional/Ricci curvature and \emph{thermodynamic length} along training or conversational trajectories; (iii) estimate multi-turn stability via \emph{holonomy gaps} and Hodge decompositions of belief flows; (iv) compare representations via \emph{Wasserstein geodesics} and \emph{persistent homology}; (v) ground early-warning signals in spectral statistics (RMT) and information-theoretic bottlenecks. Our goal is an \emph{invariant, reproducible fingerprint} $\Phi$ with confidence intervals and cost models.
  \item \emph{\textbf{nDNA v3} (dynamics, modalities, intervention).} We will (i) extend to \emph{multimodal/embodied} settings (ÆTHERs); (ii) cast alignment as a \emph{gauge field} on belief transport; (iii) formalize \emph{geometric merging} via OT barycenters; (iv) convert \emph{GENOME—VACCINE} into a \emph{control-theoretic} vaccine design that suppresses torsion spikes without degrading expressive length; (v) operationalize \emph{NEPHOS} with a \emph{Permafrost Thaw Index}—the minimal adversarial “heat” needed to activate dormant seeds—and detoxify via curvature-smoothing (Ricci-like) flows.
\end{enumerate}

\paragraph{Falsifiability, ablations, and predictive checkpoints.}
We will \emph{design to be wrong}. Concretely:
\begin{itemize}
  \item \textbf{Falsifiable links.} Pre-register hypotheses: (H1) persistent negative Ricci along alignment schedules increases multi-turn collapse risk; (H2) curvature concentration predicts jailbreak success; (H3) thermodynamic length loss predicts shortcut learning (SCAR). We will accept refutation and adjust the fingerprint.
  \item \textbf{Ablations.} (i) perturb data curricula and measure $\Delta\Phi$; (ii) alter optimizer noise/temperature and track Lyapunov/Koopman stability; (iii) inject calibrated leakage/poison patterns and test NEPHOS sensitivity vs.\ false positives.
  \item \textbf{Cross-lab replication.} Release seeds, codepaths, and audit trails to ensure the same $\Phi$ emerges across institutions and hardware stacks.
\end{itemize}

\paragraph{Operationalization: pipelines, budgets, and governance.}
We will make \emph{nDNA usable} by non-authors:
\begin{itemize}
  \item \textbf{Pipelines.} CI-ready evaluators that compute $\Phi$ nightly for tracked models; \emph{delta dashboards} highlighting where belief geometry drifts across releases.
  \item \textbf{Geometry budgets.} Product teams will set guardrails (e.g., max torsion per turn, min length under quantization, curvature bounds under alignment). Builds failing budgets trigger remediation playbooks.
  \item \textbf{Governance hooks.} Export \emph{nDNA audit cards} with versioned fingerprints, lineage maps, and risk tags (SCAR, NEPHOS, nVIRAL) for compliance, procurement, and incident response.
\end{itemize}

\paragraph{From diagnostics to interventions: vaccines, detox, and repair.}
We will not stop at measurement:
\begin{itemize}
  \item \textbf{Cognitive vaccines (GENOME—VACCINE).} Train small controllers that \emph{regularize belief transport} (reduce holonomy/torsion) while preserving semantic length; evaluate on adversarial suites and long-horizon dialogue.
  \item \textbf{Detox flows for poisoning (NEPHOS).} Identify poisoned subspaces via topological/spectral anomalies; apply \emph{targeted} curvature smoothing and distributional OT realignment; verify with \emph{Permafrost Thaw Index} reductions.
  \item \textbf{Safe merging via OT barycenters.} Replace weight-arithmetic mergers with barycentric fusions on representation laws; certify monotone risk under bounded curvature.
\end{itemize}







\section{Listen to LeCun: The Deliberate Push for ``Deep Learning''}

\noindent \textbf{Why this matters.}
In the early 2000s, ``neural nets'' carried decades of baggage. To revive representation learning and reset perceptions, a small group led by Yann LeCun, Geoffrey Hinton, and Yoshua Bengio made a \emph{deliberate rhetorical and scientific pivot}: they rallied the community around the banner of \textbf{deep learning} and backed it with concrete algorithmic results.

\begin{quote}
\emph{``Around 2003, Geoff Hinton, Yoshua Bengio and myself initiated a kind of `conspiracy' to revive the interest of the machine learning community in the problem of learning representations.''}---Y.~LeCun (2014) \cite{LeCunKDnuggets2014}
\end{quote}

\noindent \textbf{Context in three beats.}
\begin{itemize}
  \item \textbf{2003--2005: Reframe the conversation.} The trio promotes \emph{deep} representation learning (beyond 2–3 nonlinear layers), using the term ``deep learning'' as an organizing banner \cite{LeCunKDnuggets2014}.
  \item \textbf{2006--2007: Show it works.} Breakthroughs in layerwise pretraining and probabilistic generative stacks (deep belief nets) provide the first broadly convincing training schemes for deeper models \cite{Hinton2006DBN,HintonSalakhutdinov2006Science}.
  \item \textbf{2009: Codify the agenda.} Bengio’s monograph consolidates the case for \emph{learning deep architectures}, connecting unsupervised pretraining, optimization, and generalization \cite{Bengio2009LDA}.
  \item \textbf{2015: Consolidation.} A field-defining Nature review (LeCun, Bengio, Hinton) formalizes the scope, successes, and open problems of deep learning \cite{LeCun2015Nature}.
\end{itemize}

\noindent \textbf{Why wording mattered (and how it relates to \textit{nDNA}).}
``Deep learning'' was not a post-hoc label; it was a \emph{purposeful reframing} that shifted perceptions and research priorities while new results made the term stick. Our use of \textbf{Neural DNA (nDNA)} follows the same playbook: a clear name, a unifying lens (the geometry of latent belief), and a growing body of math and experiments. Naming organized the deep-learning renaissance; we expect precise definitions and reproducible results to do the same for \emph{nDNA}.

%  end introduction

\clearpage
\newpage


% \newpage

\bibliographystyle{acl_natbib}
\bibliography{custom}

\clearpage
\newpage

\noindent 


% Requires: \usepackage{graphicx}
% Optional (better float placement in two-column): \usepackage{dblfloatfix}

\newlength{\panelht}
\setlength{\panelht}{0.62\textheight} % ← overall column height (tune as you like)





\begin{figure*}[ht!]
  \centering

  % Keep both panels on one line
  \noindent\makebox[\textwidth][c]{%
    % ---------- LEFT: single tall image ----------
    \begin{minipage}[t]{0.49\textwidth}\vspace{0pt}\centering
      \includegraphics[width=\linewidth,height=\panelht,keepaspectratio]{pipe/nano_gpt.png}
      \par\vspace{0.45em}\small
      \textbf{nano-gpt (\emph{structure}).} \emph{Architecture as channel blueprint:} depth acts like the axial coordinate; \textbf{residuals} $\!\leftrightarrow\!$ \emph{bypass pipes}; \textbf{attention} / \textbf{MLP} blocks act as \emph{mixers/valves} that locally reshape the flow of representations.
    \end{minipage}\hfill
    % ---------- RIGHT: two stacked images (total height = \panelht) ----------
    \begin{minipage}[t]{0.49\textwidth}\vspace{0pt}\centering
      \includegraphics[width=\linewidth,height=.49\panelht,keepaspectratio]{pipe/flow_simulation.png}
      \par\vspace{0.3em}\small
      \textbf{Flow simulation (\emph{analogue}).} \emph{Fluid:} colored streamlines show speed through a bend and throat—\textbf{curvature} rises, \textbf{shear} increases, small \emph{recirculation} pockets may form. \emph{Semantic:} bends $\Rightarrow$ \textbf{spectral curvature} spikes ($\kappa$); constrictions $\Rightarrow$ \textbf{thermodynamic length} bursts ($\Delta L$); eddies $\Rightarrow$ local rotation in the \textbf{belief field} ($\nabla\!\times\!\mathbf{v}$).
      \par\vspace{0.75em}
      \includegraphics[width=\linewidth,height=.49\panelht,keepaspectratio]{pipe/pipe.jpg}
      \par\vspace{0.3em}\small
      \textbf{Pipeline metaphor (\emph{macro view}).} \emph{Geometry governs transport:} routing capacity and effort depend on the network of ducts. \emph{Semantic:} model design / fine-tuning shapes \textbf{where meaning flows easily}, \textbf{where it pays}, and \textbf{where it recirculates}.
    \end{minipage}%
  }

  \caption{\textbf{Semantic hydrodynamics.}
  \textbf{\emph{Model.}} We read the forward pass as \emph{semantic hydrodynamics}: a prompt injects \emph{semantic mass} that is transported through depth like a fluid through a shaped channel.
  \textbf{\emph{Why.}} Weight/attention coordinates can change without altering behavior; the \emph{on-input flow} provides \textbf{behavior-first}, \textbf{coordinate-free} signals.
  \textbf{\emph{Reading guide.}} \textbf{Bend} $\to$ \emph{spectral curvature} $\kappa$ (sharp reroutes vs.\ laminar refinement); \textbf{Pay} $\to$ \emph{thermodynamic length} $L$ (where the model expends effort; $\Delta L$ bursts mark \emph{bottlenecks}); \textbf{Push} $\to$ \emph{belief field} $\mathbf{v}$ (direction/magnitude of local drive; eddies indicate \emph{recirculation}).
  \textbf{\emph{Benefit.}} The same metaphor specifies \textbf{where to measure}—\emph{bends}, \emph{throats}, and \emph{eddies}—turning inner computation into \textbf{actionable diagnostics} and \textbf{governance thresholds}.}
  \label{fig:semantic-hydrodynamics-overview}
\end{figure*}





\noindent
\section{\textbf{Rationale and Formalization:} Why trajectories, not weights: the case for \textbf{nDNA}}

\noindent
The usual levers for interpreting and governing LLMs---parameter counts, sparsity patterns, attention heatmaps---live in coordinates that are \emph{non-identifiable} and only weakly tethered to deployed behavior.
Permutations, rotations, and low-rank re-expressions can leave the realized function intact while scrambling weight-level narratives \citep{garipov2018loss,draxler2018essentially,li2018visualizing,entezari2022role,ainsworth2023git,wortsman2022model}.
Attention visualizations, while illuminating, are not guaranteed to be \emph{faithful} causal mechanisms and drift across heads/checkpoints \citep{jain2019attention,wiegreffe2019attention,serrano2019attention,clark2019bert,michel2019sixteen,voita2019analyzing}.
By contrast, what remains stable under such reparameterizations is the \emph{on-input computation}: for a prompt $x$, the forward pass traces a \textbf{trajectory of hidden states} through depth.
Endowing representation space with information geometry (e.g., \textbf{Fisher--Rao} pullbacks) yields \textbf{coordinate-free} notions of distance, bending, and effort that track changes in the output law \citep{efron1975defining,amari2000methods,amari2016information}.
We read this as \textbf{semantic hydrodynamics}: \textbf{meaning is transported} through layers like a fluid through a shaped conduit.


\noindent
\subsection{Limits of weight–space and attention views}
Weight–space indicators (parameter counts, sparsity, individual neurons/heads) live in \emph{non-identifiable} coordinates: permutations, rotations, or refactorings can leave behavior unchanged while rewriting any weight-level narrative. Attention maps are largely \emph{descriptive}, not reliably causal or stable—different patterns can yield the same outputs and head roles drift across training. These limits motivate a \textbf{behavior-first}, \textbf{coordinate-free} view that reads the model’s \emph{on-input trajectory} of representations, rather than static weights or raw attention.


\begin{itemize}[leftmargin=*,itemsep=0.55em]

  \item \textbf{Weight space is non-identifiable and behavior-misaligned.}
  \emph{Permutation symmetries, rotations, and low-rank re-expressions} can preserve the function while scrambling weight-level narratives. Empirically, independently trained solutions are often \emph{mode-connected} by low-loss paths or become connected after accounting for permutations, undermining explanations that cling to specific coordinates \citep{garipov2018loss,draxler2018essentially,li2018visualizing,entezari2022role,ainsworth2023git}. Moreover, practical levers like \textbf{weight averaging/model soups} alter parameters while leaving deployed behavior similar or improved, again decoupling “where weights sit’’ from \emph{what the model does} \citep{wortsman2022model}. In short, \textbf{we deploy behaviors, not weights}; coordinate-specific stories are fragile.

  \item \textbf{Attention is informative but not a faithful, stable mechanism by itself.}
  Extensive tests show that \emph{similar outputs can arise from disparate attention patterns}, and directly perturbing attention often leaves predictions largely unchanged; hence attention weights are, at best, \emph{descriptive} \citep{jain2019attention,serrano2019attention}. Redundancy and role-drift are common: many heads can be pruned with little loss, a few heads do the “heavy lifting,” and head functions shift across training or fine-tuning, weakening governance value of raw maps \citep{michel2019sixteen,voita2019analyzing,clark2019bert,kovaleva2019revealing}. Post-hoc corrections (e.g., \emph{attention flow/rollout}) improve alignment with token importance but still treat attention as \emph{signals}, not ground-truth causes \citep{abnar2020quantifying}.
  Beyond attention, \textbf{critical computation lives in MLPs}: feed-forward layers behave like \emph{key–value memories} that store and retrieve factual associations, so attention alone under-specifies mechanism \citep{geva2021ffkv}.
  Methodologically, the broader saliency literature warns that visually plausible explanations can fail \emph{sanity checks}, and “faithfulness’’ must be defined and evaluated explicitly \citep{adebayo2018sanity,jacovi2020faithfulness}.

  \item \textbf{What \emph{is} stable: the on-input trajectory.}
  For each prompt \(x\), the forward pass traces a depth-indexed path of hidden states—the operational object we actually deploy. Prior analyses show that linguistic competencies emerge layerwise in consistent \emph{pipelines} (POS $\rightarrow$ parsing $\rightarrow$ NER $\rightarrow$ SRL $\rightarrow$ coreference), supporting the intuition that the \emph{trajectory through representation space} is a robust behavioral signature \citep{tenney2019bert,clark2019bert}. This motivates nDNA’s choice to work in \textbf{trajectory space}, not parameter space.

\end{itemize}


\medskip
\noindent
\subsection{Why semantic hydrodynamics matters - (deeper intuition)}
\begin{itemize}[leftmargin=*,itemsep=0.4em]
  \item \textbf{We govern \emph{behavior}, not coordinates.} Operational concerns---\emph{robustness, safety, bias, faithfulness}---attach to what the model \textbf{does} on an input, not to how its weights are labeled. Two checkpoints can behave the same while their parameters and attention differ. In short: \textbf{the weights are the map; the trajectory is the territory}.
  \item \textbf{Invariance beats introspection.} Coordinate-bound stories change under neuron permutations, subspace rotations, or low-rank refactorings; the \textbf{path an input carves} and its \textbf{geometry} (length, curvature, alignment) are \textbf{invariant} because they are measured by \textbf{how predictions would change}, not by which index moved.
  \item \textbf{Geometry turns cognition into observables.} An information metric acts as local \textbf{stiffness}: soft directions barely affect the output; stiff directions swing the predictive law. With that ruler, we quantify \textbf{how far} the model travels to reshape belief (thermodynamic length $L$), \textbf{where} it turns its internal argument (spectral curvature $\kappa$), and \textbf{what} pushes change locally (belief field $\mathbf{v}$) \citep{sivak2012thermodynamic,hyvarinen2005estimation}.
  \item \textbf{The hydrodynamics metaphor is operational.} Like fluid in a channel, semantic flow shows \textbf{corners, constrictions, and eddies}: sharp bends $\Rightarrow$ high $\kappa$; narrow throats $\Rightarrow$ bursts in $\Delta L$; local recirculation $\Rightarrow$ rotational structure in $\mathbf{v}$. These are \textbf{measurable}, per-layer signals on the actual computation.
\end{itemize}

\medskip
\noindent\textbf{\emph{What this buys us} (concrete payoffs).}
\begin{itemize}[leftmargin=*,itemsep=0.4em]
  \item \textbf{Behavior-first invariance.} Reading $\kappa$, $L$, and $\mathbf{v}$ on the trajectory yields \textbf{fingerprints} that are \textbf{comparable} across models, seeds, and checkpoints---even when weights or head roles reshuffle.
  \item \textbf{Local diagnostics.} \textbf{$\kappa$ spikes} flag brittle decision pivots; \textbf{$\Delta L$ bursts} expose capacity bottlenecks or lossy transformations; \textbf{low alignment} (small $\cos\theta$ between $\mathbf{v}$ and the tangent $\mathbf{T}$) marks layers that \textbf{move without updating belief} (staging or detours).
  \item \textbf{Governance hooks.} \textbf{Geometry budgets and thresholds}---max $\kappa$, allowable $\Delta L$ per slice, minimum alignment---become \textbf{pre-release gates}; nDNA fingerprints support \textbf{drift monitoring} after fine-tuning, pruning, quantization, or alignment.
  \item \textbf{Comparative forensics.} Because $\kappa/L/\mathbf{v}$ are tied to the output law, we can \textbf{attribute performance deltas} to \textbf{where in depth} the flow changed (e.g., a new bend from fine-tuning, an effort spike from quantization) instead of to unstable weight indices.
\end{itemize}

\medskip
\noindent\textbf{\emph{Rule of thumb}.}
If the goal is to \textbf{explain}, \textbf{compare}, or \textbf{govern} deployed behavior, analyze the \textbf{flow of meaning} that the input actually experiences. In nDNA: \textbf{curvature} says \emph{where it bends}, \textbf{thermodynamic length} says \emph{how much it pays}, and the \textbf{belief field} says \emph{what pushes it}---all with a ruler calibrated to the model's own predictions.









% Requires: \usepackage{graphicx} \usepackage{amsmath}
%\newlength{\panelht}
\setlength{\panelht}{0.27\textheight} % ← tune panel height so two rows fit nicely on the page

\begin{figure*}[t]
  \centering

  %======================== ROW 1 ========================%
  \noindent\makebox[\textwidth][c]{%
    % A — Laminar flow (fluid first, then semantic parallel)
    \begin{minipage}[t]{0.49\textwidth}\centering
      \includegraphics[width=\linewidth,height=\panelht,keepaspectratio]{pipe/laminar_flow.png}
      \par\vspace{0.4em}\small
      \textbf{Laminar flow.} \emph{Fluid:} viscous–dominated, low–Re regime; nearly parallel streamlines, negligible cross–stream mixing, no recirculation.\\
      \textbf{LLM Semantic Flow:} uniformly low spectral curvature $\kappa$, small steady $\Delta L$, and high alignment between the step and the belief push (steady refinement).
    \end{minipage}\hfill
    % B — Spectral curvature (fluid bend -> kappa spike)
    \begin{minipage}[t]{0.49\textwidth}\centering
      \includegraphics[width=\linewidth,height=\panelht,keepaspectratio]{pipe/spectral_flow.png}
      \par\vspace{0.4em}\small
      \textbf{Spectral curvature ($\kappa$) — turbulent.} \emph{Fluid:} a bend induces sharp turning, higher shear, possible separation. \\
      \textbf{LLM Semantic Flow:} a localized $\kappa$ spike at the turning point marks a sharp reroute in representation space; quasi–linear segments before/after indicate a discrete semantic pivot (e.g., topic jump, shortcut, policy jolt).
    \end{minipage}%
  }

  \vspace{0.75em}

  %======================== ROW 2 ========================%
  \noindent\makebox[\textwidth][c]{%
    % C — Thermodynamic length (fluid dissipation -> semantic effort)
    \begin{minipage}[t]{0.49\textwidth}\centering
      \includegraphics[width=\linewidth,height=\panelht,keepaspectratio]{pipe/thermodynamics_flow.png}
      \par\vspace{0.4em}\small
      \textbf{Thermodynamic length ($L$).} \emph{Fluid:} a constriction raises shear and pressure drop; energy dissipates fastest in the throat. \\
      \textbf{LLM Semantic Flow:} a stiffer metric band (hatched) and a rise in $\Delta L$ reveal a bottleneck where extra \emph{semantic effort} is paid to reshape belief (friction, detours, boundary crossing).
    \end{minipage}\hfill
    % D — Belief field (fluid velocity field -> semantic push)
    \begin{minipage}[t]{0.49\textwidth}\centering
      \includegraphics[width=\linewidth,height=\panelht,keepaspectratio]{pipe/belief_vector_flow.png}
      \par\vspace{0.4em}\small
      \textbf{Belief field ($\mathbf{v}$).} \emph{Fluid:} the velocity field sets transport; eddies (local curl) mark recirculation; alignment with streamlines indicates efficient conveyance. \\
      \textbf{LLM Semantic Flow:} $\mathbf{v}$ is the local push that most steeply changes the output law; longer arrows $\Rightarrow$ larger $\|\mathbf{v}\|$, and the side gauge shows $\cos\theta$ between $\mathbf{v}$ and the path tangent $\mathbf{T}$; circular loops on waves depict local recirculation that can trap or reinforce beliefs.
    \end{minipage}%
  }

  \caption{\textbf{LLM as an input$\to$output semantic channel.}
\emph{Model:} we read the forward pass as \emph{semantic hydrodynamics}—a prompt injects semantic mass that is transported through depth like a fluid through a shaped conduit.
\textbf{Bend} (\emph{top row}): curvature $\kappa$ distinguishes \emph{laminar} refinement from \emph{sharp} reroutes.
\textbf{Pay} (\emph{bottom left}): thermodynamic length $L$ localizes where effort concentrates via $\Delta L$ bursts (\emph{bottlenecks}).
\textbf{Push} (\emph{bottom right}): the belief field $\mathbf{v}$ reveals whether a layer update directly \emph{advances belief} (high alignment) or \emph{reorganizes information} (low alignment); eddies signal \emph{local recirculation}.\\
\textbf{Why this lens:} weight--space and attention views are \emph{non--identifiable} and unstable across checkpoints; nDNA instead reads the \emph{on--input trajectory} and its information geometry, yielding \emph{coordinate--free}, behavior--first measurements.\\
\textbf{Vision:} treat inner computation as a \emph{measurable flow} so that bends, effort, and push become quantifiable traits of cognition—comparable across inputs, layers, models, and training phases.\\
\textbf{Benefits:} \emph{actionable diagnostics}—$\kappa$ spikes flag brittle turns, $\Delta L$ bursts expose capacity bottlenecks, low $\cos\theta$ (between $\mathbf{v}$ and the tangent $\mathbf{T}$) indicates movement that does not immediately update belief; \emph{stable comparability}—geometry--based fingerprints are robust to neuron permutations and head--role drift; \emph{governance hooks}—set thresholds on $\kappa$ or $\Delta L$, track fingerprint drift after fine--tuning/pruning, and audit capacity before release.}
  \label{fig:ndna_semantic_flow_minipage}
\end{figure*}


















\end{document}
